\documentclass[10pt,a4paper]{article}
\usepackage{amsmath,amssymb,amsthm,graphicx}
% Card-local minimal preamble (journal unit sigma-2024-003-optimal-transport-and-generalized-ricci-flow, P2-ahl16 precedent).
%% Verbatim from the pinned source: the author theorem environments and macro block,
%% so every macro and environment of the counted spans is available.
% Omitted and why: the documentclass line (the journal class sigma is not on the host;
%% the assembler supplies the class) and page-geometry/font blocks (they fight the
%% forced xelatex engine).
\usepackage{amsmath,amssymb,amsthm,mathtools}

\numberwithin{equation}{section}
\newtheorem{thm}{Theorem}[section]
\newtheorem{cor}[thm]{Corollary}
\newtheorem{Lemma}[thm]{Lemma}
\newtheorem{prop}[thm]{Proposition}
 { \theoremstyle{definition}
\newtheorem{defn}[thm]{Definition}
\newtheorem{Note}[thm]{Note}
\newtheorem{Example}[thm]{Example}
\newtheorem{rmk}[thm]{Remark} }
\newcommand{\del}{\partial}
\newcommand{\delb}{\bar{\partial}}\newcommand{\dt}{\frac{\partial}{\partial t}}
\newcommand{\brs}[1]{\left| #1 \right|}
\newcommand{\G}{\Gamma}
\newcommand{\gG}{\Gamma}
\renewcommand{\gg}{\gamma}
\newcommand{\eps}{\epsilon}
\newcommand{\gD}{\Delta}
\newcommand{\gd}{\delta}
\newcommand{\gs}{\sigma}
\newcommand{\gS}{\Sigma}
\newcommand{\gu}{\upsilon}
\newcommand{\gU}{\Upsilon}
\newcommand{\gk}{\kappa}
\newcommand{\gl}{\lambda}
\newcommand{\gt}{\theta}
\newcommand{\gw}{\omega}
\newcommand{\ga}{\alpha}
\newcommand{\gb}{\beta}
\newcommand{\gL}{\Lambda}
\newcommand{\N}{\nabla}
\newcommand{\FF}{\mathcal F}
\newcommand{\WW}{\mathcal W}
\newcommand{\DD}{\mathcal D}
\newcommand{\JJ}{\mathcal J}
\renewcommand{\SS}{\mathcal S}
\newcommand{\EE}{\mathcal E}
\newcommand{\LL}{\mathcal L}
\newcommand{\CC}{\mathcal C}
\newcommand{\R}{\mathbb R}
\newcommand{\til}[1]{\widetilde{#1}}
\newcommand{\nm}[2]{\brs{\brs{ #1}}_{#2}}
\newcommand{\gf}{E}
\newcommand{\boxit}[1]{\fbox{\mbox{$#1$}}}
\newcommand{\ohat}[1]{\overset{\circ}{#1}}
\newcommand{\KN}{{\mathbin{\bigcirc\mkern-15mu\wedge}}}
\renewcommand{\bar}[1]{\overline{#1}}
\newcommand{\dtau}{\frac{d}{d \tau}}
\renewcommand{\i}{\sqrt{-1}}
\renewcommand{\part}{\del}
\newcommand{\bpart}{\bar{\del}}
\newcommand{\bbJ}{\mathbb J}
\newcommand{\ba}{\bar{a}}
\newcommand{\bb}{\bar{b}}
\newcommand{\bc}{\bar{c}}
\newcommand{\bd}{\bar{d}}
\newcommand{\be}{\bar{e}}
\newcommand{\bg}{\bar{g}}
\newcommand{\bh}{\bar{h}}
\newcommand{\bi}{\bar{i}}
\newcommand{\bj}{\bar{j}}
\newcommand{\bk}{\bar{k}}
\newcommand{\bl}{\bar{l}}
\newcommand{\bm}{\bar{m}}
\newcommand{\bn}{\bar{n}}
\newcommand{\bp}{\bar{p}}
\newcommand{\bq}{\bar{q}}
\newcommand{\br}{\bar{r}}
\newcommand{\bs}{\bar{s}}
\newcommand{\bu}{\bar{u}}
\newcommand{\bv}{\bar{v}}
\newcommand{\bw}{\bar{w}}
\newcommand{\bx}{\bar{x}}
\newcommand{\by}{\bar{y}}
\newcommand{\bz}{\bar{z}}
\newcommand{\IP}[1]{\left<#1\right>}
\newcommand{\wnorm}[1]{\left|#1\right|^2_{\frac1k}}
\newcommand{\wwnorm}[1]{\left|#1\right|^2_{\frac{k-1}k}}
\newcommand{\dd}[1]{\frac{\del}{\del #1}}
\newcommand{\bga}{\bar{\alpha}}
\newcommand{\bgb}{\bar{\beta}}
\newcommand{\bgg}{\bar{\gamma}}
\newcommand{\bmu}{\bar{\mu}}
\newcommand{\bnu}{\bar{\nu}}
\newcommand{\bgs}{\bar{\sigma}}
\DeclareMathOperator{\Var}{Var}
\DeclareMathOperator{\USC}{USC}
\DeclareMathOperator{\Sym}{Sym}
\DeclareMathOperator{\Rc}{Rc}
\DeclareMathOperator{\Rm}{Rm}
\DeclareMathOperator{\inj}{inj}
\DeclareMathOperator{\tr}{tr}
\DeclareMathOperator{\Ker}{Ker}
\DeclareMathOperator{\Img}{Im}
\DeclareMathOperator{\Id}{Id}
\DeclareMathOperator{\divg}{div}
\DeclareMathOperator{\grad}{grad}
\DeclareMathOperator{\Vol}{Vol}
\DeclareMathOperator{\diam}{diam}
\DeclareMathOperator{\im}{Im}
\DeclareMathOperator{\Area}{Area}
\DeclareMathOperator{\Ca}{Ca}
%\DeclareMathOperator*{\argmax}{argmax}
%\DeclareMathOperator{\supp}{supp}
\DeclareMathOperator{\Length}{Length}
\DeclareMathOperator{\I}{I}
\DeclareMathOperator{\End}{End}
\DeclareMathOperator{\rank}{Rank}
\DeclareMathOperator{\Ann}{Ann}
\DeclareMathOperator{\Symp}{Symp}
\DeclareMathOperator{\Ham}{Ham}
\DeclareMathOperator{\Diff}{Diff}
\DeclareMathOperator{\Vect}{Vect}
\usepackage{cleveref}
\crefname{originalcorollary}{original corollary}{original corollarys}
\Crefname{originalcorollary}{Original corollary}{Original corollarys}
\makeatletter
\newlabel{cor: cost}{{4.6}{source}{}{sourcepointer-1}{}}
\newlabel{cor: cost@cref}{{[originalcorollary][0][]4.6}{[1][0][]source}{}{}{}}
\makeatother
\crefname{originalequation}{original equation}{original equations}
\Crefname{originalequation}{Original equation}{Original equations}
\makeatletter
\newlabel{eq: heatdensity}{{3.3}{source}{}{sourcepointer-2}{}}
\newlabel{eq: heatdensity@cref}{{[originalequation][0][]3.3}{[1][0][]source}{}{}{}}
\makeatother
\crefname{originalequation}{original equation}{original equations}
\Crefname{originalequation}{Original equation}{Original equations}
\makeatletter
\newlabel{eq: heatmeasure}{{3.2}{source}{}{sourcepointer-3}{}}
\newlabel{eq: heatmeasure@cref}{{[originalequation][0][]3.2}{[1][0][]source}{}{}{}}
\makeatother
\crefname{originallemma}{original lemma}{original lemmas}
\Crefname{originallemma}{Original lemma}{Original lemmas}
\makeatletter
\newlabel{l:rhoBHE}{{3.5}{source}{}{sourcepointer-4}{}}
\newlabel{l:rhoBHE@cref}{{[originallemma][0][]3.5}{[1][0][]source}{}{}{}}
\makeatother
\title{Monotonicity of the entropy-type energy functional along the backward heat equation}
\author{Kopfer, Eva, Streets, Jeffrey}
\date{doi:10.3842/SIGMA.2024.003; CC BY 4.0}
\begin{document}
\maketitle
\setcounter{section}{4}
\setcounter{thm}{6}
\noindent\textit{Editorial note.} All delivered mathematical text is the author's own, copied from the pinned journal source file sources/sigma-2024-003-optimal-transport-and-generalized-ricci-flow/pinned-original-source.tex (Optimal Transport and Generalized Ricci Flow); no statement or proof step was rewritten, reordered or removed.\par\medskip
\noindent\textit{Editorial note.} The counted proof is the author's complete argument: the census flags the proof block as shared between cor: cost (6.2) and this corollary (6.3); a direct read of the proof resolves the pairing - the proof cites cor: cost as an input at proof line 699 and concludes the monotonicity of F claimed here, so the proof belongs to this corollary and cor: cost carries no dedicated proof of its own; no proof step is omitted.\par\medskip
\noindent\textit{Editorial note.} The label c:Fmonotone names the counted statement, and the proof environment immediately following it in the source is delivered verbatim as its proof.\par\medskip
\noindent\textit{Editorial note.} The counted statement is carried by the author's own cor environment in the article source.\par\medskip
\noindent\textit{Editorial note.} Journal-unit assembly: card-local preamble (author theorem environments and macros verbatim), article class, no bibliography in the counted span (the span carries no \textbackslash{}cite commands); references leaving the span are declared as navigation stubs carrying the article's printed numbers.\par\medskip
\noindent\textit{Editorial scope.} the counted statement and its own proof environment only\par\medskip
\noindent\textit{Allowed standard prerequisites (editorial).}\begin{itemize}
\item The backward heat equation in measure form (eq: heatmeasure). A smooth one-parameter family of probability measures (mu\_t) under generalized Ricci flow satisfies the backward heat flow when d\_t mu = -Delta mu.. Cited at statement line 691 (the hypothesis: mu solves the backward heat equation) and again at proof line 699 through the shifted endpoints mu(t'+u), mu(t''+u); declared as a navigation stub carrying the article's printed number.
\item The backward heat equation in density form (eq: heatdensity). For mu\_t = rho\_t e\textasciicircum{}\{-f\_t\} dV\_t the backward heat flow is equivalent to d\_t rho = -Delta\_f rho + R\textasciicircum{}\{H,f\} rho on the density.. Cited at proof line 708 to substitute d\_t rho and identify grad phi with grad log rho; declared as a navigation stub.
\item The backward-heat-equation lemma (l:rhoBHE, printed 3.5). The equivalence lemma between the measure form and the density form of the backward heat equation.. Invoked at proof line 699 together with Corollary cor: cost to transfer the cost monotonicity to the shifted endpoints; declared as a navigation stub.
\item The cost monotonicity corollary (cor: cost, printed 6.2). Under the hypothesis of Proposition prop: cost, if each mu\_u is a minimizer for E\_0\textasciicircum{}\{t'+u,t''+u\} with endpoints satisfying the backward heat equation, then u maps to C\_0\textasciicircum{}\{t'+u,t''+u\}(mu\_u(t'+u), mu\_u(t''+u)) nondecreasingly.. Invoked at proof line 699 to bound the shifted cost quotient by the unshifted one; declared as a navigation stub.
\end{itemize}\medskip
\ifdefined\hypertarget\hypertarget{sourcepointer-1}{}\fi
\noindent\textit{Source reference.} Original corollary 4.6: sources/sigma-2024-003-optimal-transport-and-generalized-ricci-flow/pinned-original-source.tex, source lines 683--683. Navigation pointer invoked at line(s) 699: the counted statement is phrased in notation from that statement of the same article, which is not retained in this card, so the reference is declared as a pointer carrying the number the article itself prints for it (the section and the place of the statement inside that section, derived by the census v3.1 rule for the shared theorem counter); the cited statement is an allowed input and none of it is reproduced inside the span.\par\medskip
\ifdefined\hypertarget\hypertarget{sourcepointer-2}{}\fi
\noindent\textit{Source reference.} Original equation 3.3: sources/sigma-2024-003-optimal-transport-and-generalized-ricci-flow/pinned-original-source.tex, source lines 394--394. Navigation pointer invoked at line(s) 708: the counted text cites that numbered equation of the same article, which is not retained in this card, so the reference is declared as a pointer carrying the number the article itself prints for it (the section and the place of the equation inside that section, derived by the census v3.1 rule for the section-scoped equation counter); the cited equation is an allowed input and none of it is reproduced inside the span.\par\medskip
\ifdefined\hypertarget\hypertarget{sourcepointer-3}{}\fi
\noindent\textit{Source reference.} Original equation 3.2: sources/sigma-2024-003-optimal-transport-and-generalized-ricci-flow/pinned-original-source.tex, source lines 390--390. Navigation pointer invoked at line(s) 691: the counted text cites that numbered equation of the same article, which is not retained in this card, so the reference is declared as a pointer carrying the number the article itself prints for it (the section and the place of the equation inside that section, derived by the census v3.1 rule for the section-scoped equation counter); the cited equation is an allowed input and none of it is reproduced inside the span.\par\medskip
\ifdefined\hypertarget\hypertarget{sourcepointer-4}{}\fi
\noindent\textit{Source reference.} Original lemma 3.5: sources/sigma-2024-003-optimal-transport-and-generalized-ricci-flow/pinned-original-source.tex, source lines 389--389. Navigation pointer invoked at line(s) 699: the counted statement is phrased in notation from that statement of the same article, which is not retained in this card, so the reference is declared as a pointer carrying the number the article itself prints for it (the section and the place of the statement inside that section, derived by the census v3.1 rule for the shared theorem counter); the cited statement is an allowed input and none of it is reproduced inside the span.\par\medskip
\begin{cor} \label{c:Fmonotone} Suppose that $\mu\subset P^\infty(M)$ is a smooth solution of the backward heat equation~\eqref{eq: heatmeasure}. Then
\begin{align*}
 \mathcal F=\int_M\big[\brs{\N\log\rho}^2+R^{H,f}\big] {\rm d}\mu
\end{align*}
is nondecreasing in $t$.
\end{cor}

\begin{proof}
 Fix a time $t'$. Using ellipticity of the linearized geodesic equation and an inverse function theorem argument, for $t''$ all sufficiently close to $t'$ and $u > 0$ sufficiently small, the minimizing geodesic connecting $\mu(t' + u)$ and $\mu(t'' + u)$ is smooth. By Lemma \ref{l:rhoBHE} and Corollary \ref{cor: cost}, we have
 \begin{align*}
 \frac{C_0^{t'+u,t''+u}(\mu(t'+u),\mu''(t''+u))}{t''-t'}\geq \frac{C_0^{t',t''}(\mu(t'),\mu(t''))}{t''-t'}.
 \end{align*}
 Letting $t''\to t'$
 \begin{align*}
 \left. \frac12\int_M\big[\brs{\N\phi}^2+R^{H,f}\big] {\rm d}\mu\right|_{t'+u}\geq
 \left. \frac12\int_M\big[\brs{\N\phi}^2+R^{H,f}\big] {\rm d}\mu\right|_{t'}.
 \end{align*}
 As $\rho$ solves \eqref{eq: heatdensity} it follows that $\N \phi= \N \log \rho$, giving the claim.
\end{proof}

\end{document}
